\documentclass[11pt,a4paper]{article}

\usepackage[utf8]{inputenc}
\usepackage[T1]{fontenc}
\usepackage{lmodern}
\usepackage{amsmath,amssymb,amsfonts}
\usepackage{mathtools}
\usepackage{graphicx}
\usepackage{booktabs}
\usepackage{geometry}
\usepackage{hyperref}
\usepackage{physics}
\usepackage{bm}

\geometry{margin=2.5cm}

\hypersetup{
    colorlinks=true,
    linkcolor=blue!60!black,
    citecolor=blue!60!black,
    urlcolor=blue!60!black
}

\title{
\textbf{Paper 11 — Le Flot Dimensionnel de la Gravité BuP}\\
\large Dimension Spectrale, Dimension de Marche et l'Exposant Gravitationnel Effectif
}

\author{Farid Hamdad\\
Programme Gravité Quantique Bottom-Up}

\date{2026}

\begin{document}

\maketitle

\begin{abstract}
Nous introduisons la loi dimensionnelle qui relie la diffusion sur un graphe d'intrication à un exposant gravitationnel effectif dans le cadre de la Gravité Quantique Bottom-Up. À partir d'un graphe d'intrication \(W_{ij}\) et de son laplacien \(L_{\rm ent}\), nous définissons une dimension spectrale \(d_s\) et une dimension de marche \(d_w\). Nous montrons que l'exposant gravitationnel effectif est naturellement donné par
\[
\alpha_{\rm eff} = \frac{2d_s}{d_w} + d_w - 4.
\]
La condition \(\alpha_{\rm eff} = 1\) définit un point fixe newtonien. Pour la diffusion brownienne ordinaire, \(d_w = 2\), ce point fixe donne \(d_s = 3\). Le Paper 11 fournit donc la colonne vertébrale dimensionnelle utilisée plus tard dans le programme phénoménologique BuP, y compris les courbes de rotation des galaxies et les tests de lentillage fort.
\end{abstract}

\section{Introduction}

Le programme Gravité Quantique Bottom-Up propose que la géométrie et la gravité émergent de la structure de l'intrication dans un état quantique fini.

Les papiers précédents ont introduit le graphe d'intrication,
\[
W_{ij} = I(i : j),
\]
et le laplacien d'intrication correspondant,
\[
L_{\rm ent} = D - W.
\]

Le but du Paper 11 est de dériver l'exposant d'échelle gravitationnel effectif à partir des quantités de diffusion sur graphe.

Le résultat principal est la loi dimensionnelle
\[
\boxed{
\alpha_{\rm eff} = \frac{2d_s}{d_w} + d_w - 4
}.
\]

Cette équation relie la structure microscopique du graphe à une réponse gravitationnelle émergente.

\section{Graphe d'intrication et laplacien}

Nous considérons un graphe pondéré \(G = (V, E, W)\), où les poids encodent la connectivité par information mutuelle ou un profil d'intrication effectif à grain grossier.

Le laplacien du graphe est
\[
L_{\rm ent} = D - W,
\]
avec
\[
D_{ii} = \sum_j W_{ij}.
\]

Le graphe est interprété comme le support discret sur lequel la géométrie émergente et la diffusion sont définies.

\section{Dimension spectrale}

La dimension spectrale est définie à partir de la probabilité de retour de la chaleur.

Soit
\[
K(t) = e^{-t L_{\rm ent}}.
\]

La probabilité de retour moyenne est
\[
P(t) = \frac{1}{N} \operatorname{Tr} K(t).
\]

Si le graphe possède un régime d'échelle effectif, alors
\[
P(t) \sim t^{-d_s/2}.
\]

Par conséquent,
\[
d_s = -2 \frac{d \log P(t)}{d \log t}.
\]

Dans les applications numériques, \(d_s\) est estimé par un ajustement log-log sur une fenêtre de diffusion stable.

\section{Dimension de marche}

La dimension de marche \(d_w\) est définie à partir du déplacement quadratique moyen sur le graphe.

Si
\[
\langle r^2(t) \rangle \sim t^{2/d_w},
\]
alors
\[
d_w = \frac{2}{d \log \langle r^2(t) \rangle / d \log t}.
\]

La dimension de marche mesure la vitesse à laquelle la diffusion explore le graphe.

Pour la diffusion brownienne ordinaire,
\[
d_w = 2.
\]

Pour la diffusion anomale,
\[
d_w \neq 2.
\]

\section{Dérivation dimensionnelle de l'exposant effectif}

Le cadre BuP suppose que la réponse gravitationnelle effective est contrôlée par deux ingrédients :

\begin{enumerate}
    \item le nombre de canaux de diffusion effectifs, contrôlé par \(d_s\);
    \item l'échelle de propagation anomale, contrôlée par \(d_w\).
\end{enumerate}

L'exposant effectif s'écrit donc
\[
\alpha_{\rm eff} = \frac{2d_s}{d_w} + d_w - 4.
\]

Le premier terme,
\[
\frac{2d_s}{d_w},
\]
encode le volume dimensionnel accessible par diffusion.

Le second terme,
\[
d_w - 4,
\]
encode la correction de propagation par rapport à la référence relativiste quadridimensionnelle.

Ainsi, la réponse gravitationnelle effective n'est pas assignée arbitrairement. Elle est générée par la paire de diffusion du graphe \((d_s, d_w)\).

\section{Point fixe newtonien}

Le point fixe newtonien (ou baryonique) est défini par
\[
\alpha_{\rm eff} = 1.
\]

En utilisant la loi dimensionnelle BuP,
\[
\frac{2d_s}{d_w} + d_w - 4 = 1.
\]

Par conséquent,
\[
2d_s = d_w (5 - d_w),
\]
ou
\[
d_s = \frac{d_w (5 - d_w)}{2}.
\]

Pour la diffusion brownienne,
\[
d_w = 2,
\]
nous obtenons
\[
d_s = 3.
\]

Ainsi, le régime newtonien standard correspond au point fixe
\[
(d_s,\; d_w,\; \alpha_{\rm eff}) = (3,\; 2,\; 1).
\]

\section{Protocole numérique}

Le pipeline numérique est :

\[
W_{ij}
\;\rightarrow\;
L_{\rm ent}
\;\rightarrow\;
P(t),\; \langle r^2(t) \rangle
\;\rightarrow\;
d_s,\; d_w
\;\rightarrow\;
\alpha_{\rm eff}.
\]

Pour chaque taille de système \(N\), nous calculons :

\[
d_s(N), \qquad d_w(N), \qquad \alpha_{\rm eff}(N).
\]

Les tendances en taille finie sont ensuite utilisées pour estimer le comportement à grand \(N\).

\section{Trajectoire en taille finie}

Les résultats numériques sont résumés par les trajectoires :

\[
d_s(N), \qquad d_w(N), \qquad \alpha_{\rm eff}(N).
\]

Les figures montrent l'approche vers un régime gravitationnel effectif lorsque la taille du graphe augmente.

\begin{figure}[h]
\centering
\includegraphics[width=0.82\textwidth]{figures/fig2_ds_vs_N.png}
\caption{Dimension spectrale \(d_s\) en fonction de la taille du graphe \(N\).}
\end{figure}

\begin{figure}[h]
\centering
\includegraphics[width=0.82\textwidth]{figures/fig3_dw_vs_N.png}
\caption{Dimension de marche \(d_w\) en fonction de la taille du graphe \(N\).}
\end{figure}

\begin{figure}[h]
\centering
\includegraphics[width=0.82\textwidth]{figures/fig4_alpha_predictions_vs_N.png}
\caption{Exposant effectif \(\alpha_{\rm eff}\) prédit à partir de \(d_s\) et \(d_w\).}
\end{figure}

\section{Interprétation physique}

Le flot dimensionnel définit différents régimes :

\begin{table}[h]
\centering
\begin{tabular}{lll}
\toprule
\textbf{Régime} & \textbf{Condition} & \textbf{Interprétation} \\
\midrule
Point fixe newtonien & \(\alpha_{\rm eff} = 1\) & Comportement baryonique / newtonien standard \\
Sous-newtonien & \(\alpha_{\rm eff} < 1\) & Réponse affaiblie \\
Super-newtonien & \(\alpha_{\rm eff} > 1\) & Réponse renforcée \\
Régime de transition & \(\alpha_{\rm eff} \approx 1\) & Crossover entre phases du graphe \\
\bottomrule
\end{tabular}
\caption{Interprétation de l'exposant effectif BuP.}
\end{table}

Cette interprétation devient centrale dans les applications observationnelles ultérieures.

Dans le Paper 12, la loi dimensionnelle est utilisée pour transformer les profils de matière en corrections de courbes de rotation.

Dans le Paper 14, la même loi est utilisée pour classer les galaxies SPARC dans les régimes LOW/HIGH.

Dans le Paper 21, la condition \(\alpha_{\rm eff} \simeq 1\) réapparaît comme un point fixe de lentillage fort dans SLACS.

\section{Discussion}

Le Paper 11 établit que l'exposant gravitationnel effectif est une observable de diffusion sur graphe.

Ceci est important car l'échelle gravitationnelle n'est pas introduite comme un ajustement empirique. Elle est dérivée de la structure du graphe d'intrication par l'intermédiaire de \(d_s\) et \(d_w\).

La principale limitation est que \(d_s\) et \(d_w\) dépendent tous deux de la fenêtre de diffusion. Ce n'est pas nécessairement une faiblesse. Dans BuP, différents régimes gravitationnels peuvent correspondre à différentes échelles de propagation sur le graphe.

\section{Conclusion}

Nous avons introduit la loi du flot dimensionnel de la gravité BuP :
\[
\alpha_{\rm eff} = \frac{2d_s}{d_w} + d_w - 4.
\]

Le point fixe newtonien est défini par
\[
\alpha_{\rm eff} = 1.
\]

Pour la diffusion brownienne, cela donne
\[
d_s = 3, \qquad d_w = 2.
\]

Le Paper 11 fournit donc le pont théorique entre le graphe d'intrication microscopique et la phénoménologie gravitationnelle effective développée dans les papiers ultérieurs.

\bibliographystyle{unsrt}
\bibliography{references}

\end{document}

